\documentclass[12pt]{article}

%\textheight 9in
%\textwidth 6in
%\topmargin 0in
%\oddsidemargin 0.5in
%\evensidemargin 0.5in

\usepackage{paralist}
\title{Laboratorio 2: spaziature e stili di testo 
%e  formule matematiche pi\'u complesse
}


\begin{document}
\maketitle
\date{  }

Un documento \LaTeX {} contiene sia                           comandi che il testo vero e proprio.... tutto mischiato! 

I comandi si riconoscono perch\`e sono preceduti dal backslash $\backslash$.

Inoltre alcuni caratteri non sono direttamente visualizzabili perch\`e sono caratteri speciali di \LaTeX {} usati nei vari comandi: 
\$ \& \% \# \_ \{ \}

\bigskip %inserisce uno spazio verticale

Un documento \LaTeX {} \`e formato di paragrafi separati l'uno dall'altro tramite una o
pi\`u linee bianche:

Ecco un paragrafo molto breve.

Qui comincia un altro paragrafo. Come potete vedere,
i ritorni a capo
 non hanno
alcuna importanza 
nella
formattazione 
del paragrafo: contano solo le linee bianche. 




Tante linee bianche o una sola danno lo stesso effetto
\\ % a  capo senza indentazione - corrisponde al comando \newline.
Quest'ultimo paragrafo chiude l'esempio andando a capo senza indentazione. 
 
%alternativamente usa il comando \noindent dopo una o più linee bianche


Un paragrafo \`e composto di parole separate da segni di punteggiatura o spazi (un
ritorno a capo \`e considerato come uno spazio). Potete                   inserire                       quanti 
spazi                    volete                                tra
due                          parole (anche se ne basta uno). 

\bigskip
Il \LaTeX{} "decide" quando andare accapo. Se non ci piace il punto in cui va a capo (quelle due parole proprio non vogliamo staccarle!) allora possiamo legarle con uno "spazio indivisibile" usando una tilde.

{\it Nei fumetti di Paperino, come si chiama  il riccastro rivale  di Paperon de Paperoni?}

Meglio come: 

{\it Nei fumetti di Paperino, come si chiama  il riccastro rivale  di Paperon de~Paperoni?}

\newpage
La dimensione del testo \`e fissata all'inizio del documento. Posso poi variare la dimensione usando i comandi-dichiarazioni:
(da utilizzare dunque con delle graffe di raggruppamento)

{\tiny   tiny}
{\scriptsize  scriptsize}
{\footnotesize   footnotesize}
{\small   small }
{\normalsize  normalsize} 
{\large  large}
{\Large  Large}
{\LARGE LARGE}
{\Huge Huge}

\bigskip\bigskip\bigskip
Anche lo stile del testo \`e fissato all'inizio del documento.
\bigskip

Questa \`e una linea di testo ``scritta normale" i.e. stile roman.

{\rm Questa \`e una linea di testo ``scritta normale" i.e. stile roman.}

{\sf Questa \`e una linea di testo in stile sans serif (senza grazie)}

{\tt Questa \`e una linea di testo in stile teletype (macchina da scrivere o terminale)}


\medskip %lascia spazio verticale "medio"
Per ognuno degli stili { precedenti} poi posso cambiare il "peso" 

{\bf  e avere il grassetto  i.e. boldface.}

\medskip %lascia spazio verticale "medio"
Per ognuno degli stili precedenti poi posso cambiare la forma 

{\it e avere il corsivo i.e. italic... } 

{\sl oppure inclinato i.e. slanted...} 

{\sc o Ancora il tutto-maiuscolo i.e. small capitals}.

\medskip %lascia spazio verticale "medio"


{\it Questa \`e una linea di testo in stile italic.}

\textit{ Questa \`e una linea di testo in stile italic.}

\textbf{ Questa \`e una linea di testo in boldface}

\texttt{ Questa \`e una linea di testo in teletype}



{\em Questa \`e una linea di testo ``scritta emphasized " i.e. stile italic.}

\bigskip

Questa \`e una linea di testo {\em scritta normale} i.e. stile roman. 
 
{\it Questa \`e una linea di testo {\em  in stile italic} .}

\bigskip


{\bf \em Questa \`e una linea di testo in boldface enfatizzato}

{\huge \bf (ORRIBILE!)}.

\bigskip
NOTA: i comandi con cui definire i vari stili possono avere effetti diversi a seconda dei pacchetti caricati nel preambolo (vedi lezioni successive)
\newpage

Semplici comandi per gestire gli spazi verticali sono

\bigskip
{\tt	bigskip }

\medskip
{\tt medskip}

\smallskip
{\tt smallskip}

Comandi molto utili perch\'e producono spazio proporzionato al font-size. Poi posso anche fissare uno spazio in cm con il comando {\tt vspace	}
\vspace{5cm}

e continuare a scrivere dopo... Esiste anche un equivalente {\tt hspace	} \hspace{2cm} per lo spazio in orizzontale.

In genere si usano questi comandi in condizioni di "emergenza"... solitamente si lascia gestire la pagina a \LaTeX!

{\bf NOTA:} Piccoli spazi orizzontali (positivi o negativi) invece si usano spesso per migliorare la leggibilit\'a di formule matematiche. Per questo si usano comandi dedicati.
\newpage

Gli ambienti in \LaTeX {} sono strutture estremamente potenti di cui si fa un grosso uso.
LATEX.


 

\begin{verbatim}
\begin{ambiente}[opzione]{parametro}
 L'ambiente verbatim ignora qualunque comando \LaTeX 
e scrive in \texttt tutto         come appare 




(anche gli accapo!)
\end{ambiente}
\end{verbatim}

Esempi di ambienti:

\begin{em}
Il testo contenuto in questo ambiente \`e messo in evidenza.
\end{em}

\bigskip
Nel mezzo di un testo inserisco una citazione
\begin{quote}
In re mathematica ars proponendi quaestionem 
pluris facienda est quam solvendi.

(In mathematics the art of asking questions 
is more valuable than solving problems.)
\end{quote}
e poi continuo con il testo normale...

\bigskip

\begin{center}
questo \'e un lunghissimo testo centrato per provare l'ambiente {\it center}

articolato addirittura su due linee
\end{center}

\bigskip

\begin{flushleft}
Ecco un paragrafo giustificato a sinistra. 

se serve...

Ecco due linee \\
giustificate a sinistra.
\end{flushleft}

\bigskip
\begin{flushright}
e per non fare disparit\'a ..\\
Ecco un paragrafo 

giustificato a destra. 

\end{flushright}
\newpage

Gli ambienti per le {\LARGE liste} sono {\tt itemize, enumerate e description}.

\begin{itemize}
\item colazione 
\item pranzo
\item merenda
\item cena
\end{itemize}

e se non mi piacciono i pallini?
\begin{itemize}
\item[!] colazione
\item[?] pranzo
\item[+] cena
\end{itemize}


Posso anche innestare una lista in una lista...
\begin{itemize}
\item colazione:
	\begin{itemize}
		\item latte 
		\item caff\'e
		\item biscotti
	\end{itemize}
\item pranzo
\item merenda
\item cena
\end{itemize}
 
Oppure posso fare una lista numerata: 

\begin{enumerate}
\item colazione
\item pranzo
\item merenda
\item cena
\end{enumerate}
 
Con i numeri romani? (NOTA: devo inserire :
\begin{verbatim}
 \usepackage{paralist}
\end{verbatim})


\begin{enumerate}[(i)]
\item colazione
\item pranzo
\item merenda
\item cena
\end{enumerate}

Oppure posso fare una lista tipo dizionario:
\begin{description}
\item[colazione] si fa al mattino appena svegli
\item [pranzo] si fa a met\'a giornata
\item [merenda] nel pomeriggio... amata dai bambini!
\item [cena] a fine giornata... magari non troppo tardi!
\end{description}
\newpage
{\LARGE ... e finalmente TABELLE}

\bigskip\bigskip

\begin{tabular}{lcr}
uno &  & tre \\
quattro & cinque & sei \\
quattro & cinque & sei \\
quattro & cinque & sei \\
quattro & cinque & sei \\
sette & otto & novantanovemila
\end{tabular}
\bigskip\bigskip


\begin{tabular}{|c|c|c|}
  \hline
  % after \\: \hline or \cline{col1-col2} \cline{col3-col4} ...
  $a_{11}$ &  $a_{12}$ &  $a_{13}$ \\
  \cline{1-2}
%\hline
   $a_{21}$ &  $a_{22}$ &  $a_{23}$ \\
\hline
   $a_{31}$ &  $a_{32}$ &  $a_{33}$ \\
  \hline
\end{tabular}

\bigskip\bigskip

\begin{tabular}{||p{5cm}||*{2}{c|}|}
\hline
& Contenuto & Quantit\'a \\
\hline
\hline
 Heineken & 33 cl & 10 \\
\hline
 Guinness & 1 pinta & 5 \\
\hline
 Kronenbourg & 33 cl & 0 \\
\hline
\end{tabular} 

\bigskip\bigskip

\begin{tabular}{c r @{,} l}
%\begin{tabular}{c r l}
Espressione &
\multicolumn{2}{c}{Valore}\\
\hline
$\pi$ & 3&1416 \\
$\pi^{\pi}$ & 36&46 \\
$(\pi^{\pi})^{\pi}$ & 80662&7\\
\end{tabular}

\bigskip\bigskip

\begin{center}
\begin{tabular}{|c||r|c|}
\hline
\multicolumn{3}{|c|}{Telecom Italia}\\
\hline
\hline
&\multicolumn{2}{c|}{Prezzo}\\
\cline{2-3}
mese&min&max\\
\hline
gen 2005&\multicolumn{2}{c|}{2.238--2.274}\\
\hline
feb 2005&\multicolumn{2}{c|}{2.268--2.312}\\
\hline
mar 2005&\multicolumn{2}{c|}{2.290--2.356}\\
\hline
\end{tabular}
\end{center}

\bigskip\bigskip
Questo invece \'e un altro sistema:

\bigskip

\begin{tabbing}
  nomilunghi \= coselunghissimo \= animali \kill
  % \> for next tab, \\ for new line...
Carlo \> cartello \> cammello \\
Manfredi \> mano \> mucca \\
  Dario \> dente \> delfino 
\end{tabbing}



\end{document}

